% =====================================================================
%  2bat-ud1-16.tex  —  SESSIÓ 16: Prova escrita de la Unitat 1
%  (sense preàmbul: s'inclou des de main.tex)
%  Criteris: C1, C2, C3, C4. Model de PROVA (solucionari + barem per al professorat).
% =====================================================================
\horatitol{Sessió 16 — Prova escrita de la Unitat 1}%
{Model de prova individual amb els quatre blocs (C1–C4). Temps orientatiu: 50 minuts.}

\subsection*{Indicacions}

La prova consta de \textbf{quatre blocs}, un per cada criteri d'avaluació (C1–C4),
amb un valor de \textbf{2,5 punts} cadascun. Es valora tant la correcció del càlcul
com la \emph{comprovació de la compatibilitat}, una \emph{disposició ordenada} i la
\emph{interpretació dels resultats} dins el context. Pots usar la calculadora.
Escriu tots els passos.

% ---------------------------------------------------------------------
\subsection{Bloc 1 (C1) — Construir i llegir una matriu}

\begin{exercicis}[Bloc 1 \quad (2,5 punts)]
\begin{exlist}
  \item Una entitat anota les persones ateses en tres serveis (acollida, formació,
        inserció) durant dos trimestres. Al primer trimestre: $120$, $80$ i $50$; al
        segon: $135$, $90$ i $60$.
        \begin{enumerate}[label=\alph*), leftmargin=1.6em, topsep=2pt]
          \item Organitza les dades en una matriu (files: serveis; columnes:
                trimestres) i indica'n la dimensió.
          \item Quines persones va atendre el servei de formació el segon trimestre?
          \item Quin servei va atendre més persones en conjunt (sumant els dos
                trimestres)?
        \end{enumerate}
\end{exlist}
\end{exercicis}

\subsection{Bloc 2 (C2) — Sumar/restar i escalar}

\begin{exercicis}[Bloc 2 \quad (2,5 punts)]
\begin{exlist}
  \item El pressupost trimestral de dos serveis (files) en personal i materials
        (columnes) ve en dues partides:
        \[ A=\begin{pmatrix} 2000 & 600 \\ 1600 & 400 \end{pmatrix},\qquad
           B=\begin{pmatrix} 400 & 200 \\ 300 & 100 \end{pmatrix}. \]
        \begin{enumerate}[label=\alph*), leftmargin=1.6em, topsep=2pt]
          \item Troba la despesa total $A+B$.
          \item S'aplica una retallada lineal del $10\%$. Per quin escalar cal
                multiplicar? Troba el nou pressupost i explica què significa l'escalar.
        \end{enumerate}
\end{exlist}
\end{exercicis}

\subsection{Bloc 3 (C3) — Producte de matrices}

\begin{exercicis}[Bloc 3 \quad (2,5 punts)]
\begin{exlist}
  \item Un servei atén dos usuaris amb tres tipus de sessió:
        $S=\begin{pmatrix} 2 & 3 & 1 \\ 4 & 1 & 2 \end{pmatrix}$. Cada tipus de sessió
        dura $H=\begin{pmatrix} 1 \\ 2 \\ 3 \end{pmatrix}$ hores. Calcula $S\cdot H$ i
        digues quantes hores de dedicació genera cada usuari.
  \item \textbf{(Mecanització)} Calcula el producte
        $\begin{pmatrix} 3 & 1 \\ 2 & 4 \end{pmatrix}\cdot\begin{pmatrix} 2 & 0 \\ 1 & 5 \end{pmatrix}$.
\end{exlist}
\end{exercicis}

\subsection{Bloc 4 (C4) — Compatibilitat i interpretació}

\begin{exercicis}[Bloc 4 \quad (2,5 punts)]
\begin{exlist}
  \item Una matriu d'usuaris per serveis és $4\times 3$ i una de serveis per cost
        és $3\times 1$.
        \begin{enumerate}[label=\alph*), leftmargin=1.6em, topsep=2pt]
          \item Es pot fer el producte? Quina dimensió tindria el resultat?
          \item Què representaria cada número del resultat?
          \item Es podria fer el producte en l'ordre contrari? Justifica-ho.
        \end{enumerate}
\end{exlist}
\end{exercicis}

% ---------------------------------------------------------------------
\fullsolucions{Sessió 16}
\begin{solucions}
\textbf{Bloc 1 (C1) — 2,5 punts}
\begin{sollist}
  \item (a) $\begin{pmatrix} 120 & 135 \\ 80 & 90 \\ 50 & 60 \end{pmatrix}$,
        \;dimensió $3\times 2$. \;(b) El servei de formació (fila 2) el segon
        trimestre (columna 2): $90$ persones. \;(c) Sumant per files: acollida
        $120+135=255$, formació $80+90=170$, inserció $50+60=110$. L'acollida n'atén
        més ($255$).
\end{sollist}
\textit{Barem:} construcció correcta de la matriu i dimensió, $1$ punt; lectura de
l'element (b), $0{,}5$; conclusió justificada (c), $1$ punt.

\textbf{Bloc 2 (C2) — 2,5 punts}
\begin{sollist}
  \item (a) $A+B=\begin{pmatrix} 2400 & 800 \\ 1900 & 500 \end{pmatrix}$. \;(b) Cal
        multiplicar per l'escalar $0{,}9$ (es manté el $90\%$):
        $0{,}9\cdot(A+B)=\begin{pmatrix} 2160 & 720 \\ 1710 & 450 \end{pmatrix}$.
        L'escalar aplica el mateix percentatge de retallada a tots els conceptes
        alhora.
\end{sollist}
\textit{Barem:} suma correcta (a), $1$ punt; escalar correcte i nou pressupost (b),
$1$ punt; interpretació de l'escalar, $0{,}5$.

\textbf{Bloc 3 (C3) — 2,5 punts}
\begin{sollist}
  \item $S\cdot H=\begin{pmatrix} 2\per 1+3\per 2+1\per 3 \\ 4\per 1+1\per 2+2\per 3 \end{pmatrix}
        =\begin{pmatrix} 11 \\ 12 \end{pmatrix}$. L'usuari 1 genera $11$ hores;
        l'usuari 2, $12$ hores.
  \item $\begin{pmatrix} 3 & 1 \\ 2 & 4 \end{pmatrix}\cdot\begin{pmatrix} 2 & 0 \\ 1 & 5 \end{pmatrix}
        =\begin{pmatrix} 7 & 5 \\ 8 & 20 \end{pmatrix}$.
\end{sollist}
\textit{Barem:} producte $S\cdot H$ correcte amb interpretació, $1{,}5$ punts;
mecanització correcta, $1$ punt.

\textbf{Bloc 4 (C4) — 2,5 punts}
\begin{sollist}
  \item (a) Sí: $4\times 3$ per $3\times 1$ dona un resultat $4\times 1$. \;(b) Cada
        número és el cost total d'un usuari (files $=$ usuaris; única columna $=$
        cost). \;(c) En l'ordre contrari, $3\times 1$ per $4\times 3$, \emph{no} es
        pot fer: les columnes de la primera ($1$) no coincideixen amb les files de
        la segona ($4$).
\end{sollist}
\textit{Barem:} compatibilitat i dimensió (a), $1$ punt; interpretació (b),
$0{,}75$; ordre contrari justificat (c), $0{,}75$.

\medskip
\noindent\textbf{Nota per al professorat.} Total: $10$ punts ($2{,}5$ per bloc). A
més de la correcció del càlcul, es valora la comprovació de la compatibilitat, la
disposició ordenada i la interpretació dins el context. Aquesta prova és
l'instrument «Proves escrites» (60\% de la qualificació de la unitat). Les dades
contextuals s'han d'actualitzar amb xifres vigents en el moment d'impartir la unitat.
\end{solucions}
